\begin{comment}
  Homework solution of Calculus 2 - Chapter 4 Integral Application
  Licensed under MIT license
\end{comment}

\documentclass[12pt]{article}
\usepackage[utf8]{inputenc}
\usepackage{amsmath}
\usepackage{geometry}

\geometry{a4paper, margin=1in}

\title{No. 9 - Aplikasi Integral}
\author{Ahmad Fakhrul Bawani}
\date{}

\begin{document}

\maketitle

\section*{Solusi Volume Benda Putar}

Diketahui daerah yang dibatasi oleh kurva $y = 4\sqrt{\sin x}$, sumbu-$x$ ($y=0$), serta garis $x = \frac{\pi}{4}$ dan $x = \frac{2\pi}{3}$ diputar mengelilingi sumbu-$x$. Kita akan menghitung volumenya menggunakan metode cakram (\textit{disk method}).

\subsection*{1. Rumus Volume Metode Cakram}
Karena sumbu putar adalah sumbu-$x$, maka volume ($V$) dihitung dengan rumus:
\[ V = \int_{a}^{b} \pi [f(x)]^2 \, dx \]
Substitusi fungsi $f(x) = 4\sqrt{\sin x}$ dan batas interval $[\frac{\pi}{4}, \frac{2\pi}{3}]$:
\[ V = \pi \int_{\pi/4}^{2\pi/3} \left( 4\sqrt{\sin x} \right)^2 \, dx \]
\[ V = \pi \int_{\pi/4}^{2\pi/3} 16 \sin x \, dx \]

\subsection*{2. Menghitung Integral}
Keluarkan konstanta dan integralkan $\sin x$:
\begin{align*}
V &= 16\pi \int_{\pi/4}^{2\pi/3} \sin x \, dx \\
V &= 16\pi \left[ -\cos x \right]_{\pi/4}^{2\pi/3}
\end{align*}

Evaluasi pada batas-batas tersebut:
\begin{itemize}
  \item Untuk $x = \frac{2\pi}{3}$: $-\cos(\frac{2\pi}{3}) = -(-\frac{1}{2}) = \frac{1}{2}$
  \item Untuk $x = \frac{\pi}{4}$: $-\cos(\frac{\pi}{4}) = -\frac{\sqrt{2}}{2}$
\end{itemize}

\subsection*{3. Hasil Akhir}
\[ V = 16\pi \left( \frac{1}{2} - \left( -\frac{\sqrt{2}}{2} \right) \right) \]
\[ V = 16\pi \left( \frac{1 + \sqrt{2}}{2} \right) \]
\[ V = 8\pi (1 + \sqrt{2}) \]

\subsection*{Kesimpulan}
Jadi, volume benda putar yang dihasilkan adalah $\mathbf{8\pi (1 + \sqrt{2})}$ \textbf{satuan kubik}.

\end{document}